% 数学ノート集・共通テンプレート v2（積分論2026年度第1部の版面を参考）
% A5判。章内で定義・定理・例・注意等の番号を共通化。
% コンパイル： platex -interaction=nonstopmode math_note_template_v2.tex (2回)
%              dvipdfmx math_note_template_v2.dvi
\documentclass[a5paper,10pt,onecolumn]{jsarticle}
\usepackage[dvipdfmx,top=22mm,bottom=20mm,inner=19mm,outer=18mm,headheight=13pt,headsep=6mm]{geometry}
\usepackage{amsmath,amssymb,amsthm,mathtools}
\usepackage{enumitem,fancyhdr,titlesec}
\usepackage[dvipdfmx,bookmarks=true,colorlinks=false]{hyperref}
\usepackage{pxjahyper}

% ----- 差し替える文書情報 -----
\newcommand{\NoteTitle}{層コホモロジーの構成と基本定理}
\newcommand{\NoteSeries}{数学ノート集}

% ----- 文字・組版 -----
% 参考PDFの質素な明朝＋Computer Modernの印象を引き継ぎながら，
% A5判で窮屈にならないよう本文の行間を少し確保する．
\linespread{1.13}
\setlength{\parindent}{1zw}
\setlength{\parskip}{0pt}
\setlength{\abovedisplayskip}{6pt plus 2pt minus 2pt}
\setlength{\belowdisplayskip}{6pt plus 2pt minus 2pt}
\setlength{\abovedisplayshortskip}{3pt plus 2pt minus 1pt}
\setlength{\belowdisplayshortskip}{4pt plus 2pt minus 1pt}
\setlist[enumerate]{topsep=3pt,itemsep=1pt,parsep=0pt,leftmargin=2.4zw}
\titleformat{\section}{\normalfont\normalsize\bfseries}{\thesection}{0.8zw}{}
\titlespacing*{\section}{0pt}{2.0ex plus .4ex}{.9ex}
\titleformat{\subsection}{\normalfont\normalsize\bfseries}{\thesubsection}{0.7zw}{}
\titlespacing*{\subsection}{0pt}{1.5ex plus .3ex}{.6ex}
\raggedbottom

% ----- 柱とページ番号：ページ上部、罫線なし -----
\fancypagestyle{mathnote}{%
  \fancyhf{}%
  \fancyhead[L]{\footnotesize \NoteSeries\ （\NoteTitle）}%
  \fancyhead[R]{\footnotesize \thepage}%
  \renewcommand{\headrulewidth}{0pt}%
  \renewcommand{\footrulewidth}{0pt}%
}
\pagestyle{mathnote}

% ----- 節内通し番号の定義・定理群 -----
\newtheoremstyle{mathnotetheorem}
 {4.5pt}{4.5pt}{\normalfont}{}%
 {\bfseries}{}{0.65zw}%
 {\thmname{#1}\thmnumber{ #2}\thmnote{（#3）}}
\theoremstyle{mathnotetheorem}
\newtheorem{theorem}{定理}[section]
\newtheorem{definition}[theorem]{定義}
\newtheorem{lemma}[theorem]{補題}
\newtheorem{proposition}[theorem]{命題}
\newtheorem{corollary}[theorem]{系}
\newtheorem{example}[theorem]{例}
\newtheorem{remark}[theorem]{注意}
\renewcommand{\proofname}{証明}
\renewcommand{\qedsymbol}{$\square$}
\makeatletter
\renewenvironment{proof}[1][\proofname]{%
  \par\pushQED{\qed}\normalfont\topsep4\p@\@plus2\p@\relax
  \trivlist\item[\hskip\labelsep\bfseries #1]\ignorespaces
}{\popQED\endtrivlist\@endpefalse}
\makeatother

% ----- 使うノートで適宜拡張する記号 -----
\newcommand{\R}{\mathbb{R}}
\newcommand{\norm}[1]{\lVert #1\rVert}
\newcommand{\abs}[1]{\lvert #1\rvert}

\begin{document}
\section*{読了後にできること}
本稿を読むための前提は，層コホモロジー，acyclicな層，長完全列．
本稿の具体的な到達課題は，簡単な空間で層コホモロジーの零・非零を判定する．

\section{Abel群の層と大域切断}

本稿では位相空間 $X$ 上のAbel群の層の圏
$\operatorname{Sh}(X,\mathbf{Ab})$ を考える．
前層・層・茎・層化，および
ホモロジー代数の完全列・導来函手を前提とする．

\begin{definition}
Abel群の層 $\mathcal F$ の大域切断を
\[
 \Gamma(X,\mathcal F)=\mathcal F(X)
\]
と書く．層の射
$\varphi:\mathcal F\to\mathcal G$ に対し
$\Gamma(X,\varphi)=\varphi_X$ と置くことで，
\[
 \Gamma(X,-):\operatorname{Sh}(X,\mathbf{Ab})
 \longrightarrow\mathbf{Ab}
\]
という函手が得られる．
\end{definition}

\begin{proposition}[大域切断の左完全性]
\label{prop:global-left}
層の短完全列
\[
 0\to\mathcal F'\xrightarrow{i}\mathcal F
 \xrightarrow{p}\mathcal F''\to0
\]
に対し
\[
 0\to\Gamma(X,\mathcal F')
 \xrightarrow{i_X}\Gamma(X,\mathcal F)
 \xrightarrow{p_X}\Gamma(X,\mathcal F'')
\]
は完全である．
ただし最後の $p_X$ は一般には全射でない．
\end{proposition}
\begin{proof}
層の完全性は各茎で検査できる．
$i$ は単射な層の射なので，各開集合上の写像も
単射であり，$i_X$ は単射である．
$p_Xi_X=0$ は層の射の合成が零であることから従う．
$s\in\mathcal F(X)$ が $p_X(s)=0$ を満たすとする．
各点 $x$ では $s_x$ が
$\ker p_x=\operatorname{im}i_x$ に属するので，
ある近傍 $U_x$ で $s|_{U_x}$ は
$\mathcal F'$ の切断から来る．
この局所的な持ち上げは $i$ の単射性によって
重なり合う近傍上で一致する．
$\mathcal F'$ の貼り合わせ公理を適用して
大域切断 $s'\in\mathcal F'(X)$ を得て
$i_X(s')=s$ となる．
従って $\ker p_X=\operatorname{im}i_X$ である．
最後の全射性は，局所的な持ち上げが存在しても
それらを大域的に貼り合わせられるとは限らないため
一般には成立しない．
\end{proof}

\section{Skyscraper層と入射層}

\begin{definition}
$x\in X$ とAbel群 $A$ に対して，
\[
 (i_{x*}A)(U)=
 \begin{cases}
 A&(x\in U),\\
 0&(x\notin U)
 \end{cases}
\]
と置き，明らかな制限写像を入れる．
これを点 $x$ に支持された
\emph{skyscraper層}という．
\end{definition}

\begin{lemma}[茎函手の完全性]\label{lem:stalk-exact}
$x\in X$ について，
$\mathcal F\mapsto\mathcal F_x$ は
Abel群の層の短完全列を
Abel群の短完全列へ写す．
\end{lemma}
\begin{proof}
層の短完全性は各茎における完全性と同値である．
具体的には茎は $x$ を含む開近傍全体に沿う
切断Abel群の帰納極限として得られる．
Abel群のフィルター付き帰納極限は完全である：
核または像に属するという有限個の元に関する条件は
十分小さい共通近傍に制限すれば確認できる．
従って茎へ移した列は完全である．
\end{proof}

\begin{lemma}[skyscraper層の随伴性]\label{lem:skyscraper-adj}
任意の層 $\mathcal F$ とAbel群 $A$ に対して
\[
 \operatorname{Hom}_{\operatorname{Sh}(X)}
 (\mathcal F,i_{x*}A)
 \cong\operatorname{Hom}_{\mathbb Z}(\mathcal F_x,A)
\]
が自然に成立する．
\end{lemma}
\begin{proof}
層の射 $\varphi:\mathcal F\to i_{x*}A$ が与えられると，
$x$ の近傍での写像
$\varphi_U:\mathcal F(U)\to A$ は制限と可換なので，
同値な芽に同じ値を与える．
従って茎上の準同型
$\varphi_x:\mathcal F_x\to A$ が定まる．
逆に $\psi:\mathcal F_x\to A$ が与えられたとき，
$x\in U$ なら
$s\mapsto\psi(s_x)$，
$x\notin U$ なら零写像と定める．
これらの写像は制限と可換なので層の射となる．
両構成は互いに逆であり自然性も明らかである．
\end{proof}

\begin{theorem}[Abel群の層には十分な入射対象がある]
\label{thm:enough-injective}
任意のAbel群の層 $\mathcal F$ は
入射的な層への単射
$\mathcal F\hookrightarrow\mathcal I$ を持つ．
\end{theorem}
\begin{proof}
任意の $x\in X$ について
Abel群 $\mathcal F_x$ を入射Abel群 $E_x$ に埋め込む．
その存在はBaerの判定条件と
$\mathbb Q/\mathbb Z$ を用いた埋め込み構成から従う．
$\mathcal I_x=i_{x*}E_x$ と置くと，
補題\ref{lem:skyscraper-adj}および
補題\ref{lem:stalk-exact}によって
$\operatorname{Hom}(-,\mathcal I_x)$ は
$\operatorname{Hom}((-)_{x},E_x)$ と自然同型であり，
後者は完全函手だから $\mathcal I_x$ は入射的である．

\[
 \mathcal I=\prod_{x\in X}\mathcal I_x
\]
と置く．入射対象の直積は入射的である．
実際，短完全列にHomを適用すると直積のHomは
各成分のHomの直積となり，
Abel群の直積は全射を成分ごとに持ち上げられるので
完全性を保つ．
各切断 $s\in\mathcal F(U)$ を，
$x\in U$ のとき $s_x\in\mathcal F_x\subset E_x$ に写し，
$x\notin U$ のとき零に写すことで
層の射 $j:\mathcal F\to\mathcal I$ を定める．
この射は茎で単射である：
各 $x$ で成分 $i_{x*}E_x$ への写像を見れば
$s_x\mapsto s_x\in E_x$ となるからである．
従って $j$ は単射な層の射であり，
十分な入射対象が存在する．
\end{proof}

\section{層コホモロジーの定義}

\begin{definition}
Abel群の層 $\mathcal F$ の入射分解
\[
 0\to\mathcal F\to\mathcal I^0
 \xrightarrow{d^0}\mathcal I^1
 \xrightarrow{d^1}\cdots
\]
を取り，
\[
 H^q(X,\mathcal F)
 =H^q\bigl(\Gamma(X,\mathcal I^\bullet)\bigr)
\]
と定める．
これを第 $q$ 層コホモロジー群という．
\end{definition}

\begin{theorem}[入射分解によらない定義]\label{thm:sheaf-wd}
層コホモロジー群 $H^q(X,\mathcal F)$ は，
入射分解の選択によらず自然同型を除いて定まる．
また
\[
 H^0(X,\mathcal F)\cong\Gamma(X,\mathcal F)
\]
である．
\end{theorem}
\begin{proof}
入射分解の比較定理を用いる．
二つの入射分解 $\mathcal I^\bullet$，
$\mathcal J^\bullet$ の間に，
$\mathcal F$ の恒等射を延長するコチェイン写像が
両方向に存在し，
その合成はそれぞれ恒等写像とコチェインホモトピックである．
この比較定理は射影分解の比較定理の双対であり，
その証明は分解の各次数で入射対象の延長性を使って
帰納的に構成するものである．
大域切断函手を適用してもコチェインホモトピーは保たれるので，
二つのコホモロジーは自然同型となる．

また $H^0$ は
$\ker(\Gamma(\mathcal I^0)\to\Gamma(\mathcal I^1))$
であり，大域切断の左完全性によって
この核は $\Gamma(\mathcal F)$ と自然に同型である．
\end{proof}

\section{Flabby層と非輪状性}

\begin{definition}
Abel群の層 $\mathcal F$ が\emph{flabby}であるとは，
任意の開集合包含 $V\subset U$ について
制限写像
$\mathcal F(U)\to\mathcal F(V)$ が全射であることをいう．
\end{definition}

\begin{lemma}[flabby核の大域切断]\label{lem:flabby-lifting}
層の短完全列
\[
 0\to\mathcal A\to\mathcal B\to\mathcal C\to0
\]
で $\mathcal A$ がflabbyなら，
任意の開集合 $U$ に対して
$\mathcal B(U)\to\mathcal C(U)$ は全射である．
\end{lemma}
\begin{proof}
$c\in\mathcal C(U)$ を固定する．
層の射 $\mathcal B\to\mathcal C$ は茎で全射なので，
$U$ の各点の近傍に $c$ の局所的な持ち上げが存在する．
開集合 $V\subset U$ と
$b_V\in\mathcal B(V)$ で $b_V$ が $c|_V$ を持ち上げる対を
延長の順序で考える．
鎖の上限は層の貼り合わせによって得られるため，
Zornの補題から極大な対 $(V,b_V)$ を取れる．
もし $V\ne U$ なら，
$U\setminus V$ の点 $x$ の近傍 $W\subset U$ で
局所持ち上げ $b_W$ が存在する．
重なり $V\cap W$ では $b_V-b_W$ の像が零なので，
核 $\mathcal A$ の切断 $a$ に一意に対応する．
$\mathcal A$ はflabbyなので $a$ は $W$ の切断へ延長できる．
その像を $b_W$ から引けば $b_V$ と重なりで一致する．
貼り合わせによって $V\cup W$ 上の持ち上げが得られ，
極大性に矛盾する．
従って $V=U$ であり全射性が従う．
\end{proof}

\begin{theorem}[flabby層の非輪状性]\label{thm:flabby-acyclic}
$\mathcal F$ がflabbyなら
\[
 H^q(X,\mathcal F)=0\qquad(q\ge1).
\]
\end{theorem}
\begin{proof}
まず入射層はflabbyであることを示す．
開集合 $U\subset X$ と
入射層 $\mathcal I$ の切断 $s\in\mathcal I(U)$ を取る．
$U$ の定数Abel群層の零延長 $j_!\mathbb Z_U$ から
$\mathcal I$ への層の射は，$s$ の指定と同値である．
自然な単射
$j_!\mathbb Z_U\hookrightarrow\underline{\mathbb Z}_X$
を考える．
$\mathcal I$ の入射性によってこの射は
$\underline{\mathbb Z}_X\to\mathcal I$ へ延長でき，
その延長は $\mathcal I(X)$ の切断に対応する．
従って $s$ は大域切断へ延長できる．
同じ議論を $V\subset U$ の相対的な位相空間で行えば
任意の制限写像が全射であり，$\mathcal I$ はflabbyとなる．

flabby層 $\mathcal F$ の入射分解を
$0\to\mathcal F\to\mathcal I^0\to\mathcal I^1\to\cdots$
と取る．
$C^1=\operatorname{coker}(\mathcal F\to\mathcal I^0)$ と置く．
$\mathcal I^0$ はflabbyであり，
補題\ref{lem:flabby-lifting}により
$\mathcal I^0(U)\to C^1(U)$ は全射である．
任意の $V\subset U$ の $C^1(V)$ の切断を
$\mathcal I^0(V)$ に持ち上げ，
$\mathcal I^0$ のflabby性によって
$\mathcal I^0(U)$ へ延長してから $C^1(U)$ に写すと，
元の切断の延長を得る．
従って $C^1$ もflabbyである．
この操作を帰納的に繰り返し，
各段階の余核 $C^q$ がflabbyであることを得る．
補題\ref{lem:flabby-lifting}から
入射分解の大域切断に関する列は
各次数で完全となるので，
正の次数のコホモロジーは零である．
\end{proof}

\section{長完全コホモロジー列}

\begin{theorem}[長完全列]\label{thm:sheaf-long}
Abel群の層の短完全列
$0\to\mathcal F'\to\mathcal F\to\mathcal F''\to0$
に対し，自然な長完全列
\[
 \begin{aligned}
 0&\to H^0(X,\mathcal F')\to H^0(X,\mathcal F)
 \to H^0(X,\mathcal F'')\\
 &\to H^1(X,\mathcal F')\to H^1(X,\mathcal F)
 \to H^1(X,\mathcal F'')\to H^2(X,\mathcal F')\to\cdots
 \end{aligned}
\]
が存在する．
\end{theorem}
\begin{proof}
Abel圏の入射版horseshoe lemmaを適用する．
短完全列の両端 $\mathcal F',\mathcal F''$ に
入射分解を選ぶと，中間 $\mathcal F$ の入射分解を
各次数で直和を取る形に構成でき，
三つの入射分解の間に短完全な複体の列が得られる．
この構成は各次数で入射対象の延長性を用い，
前段階の余核の短完全列に同じ操作を繰り返す
帰納法による．
大域切断は直和を保存するため，
三つの大域切断複体の列も次数ごとに短完全である．
鎖複体の短完全列からホモロジーの長完全列を作る
連結準同型の構成を次数反転して適用すると，
表示した長完全列が得られる．
\end{proof}

\section{\v Cechコホモロジー}

\begin{definition}
開被覆 $\mathcal U=\{U_i\}_{i\in I}$ と
Abel群の層 $\mathcal F$ に対して，
$U_{i_0\cdots i_q}=U_{i_0}\cap\cdots\cap U_{i_q}$ と置き，
\[
 \check C^q(\mathcal U,\mathcal F)
 =\prod_{i_0<\cdots<i_q}
 \mathcal F(U_{i_0\cdots i_q})
\]
を\v Cechコチェイン群とする．
境界写像を
\[
 (\delta c)_{i_0\cdots i_{q+1}}
 =\sum_{j=0}^{q+1}(-1)^j
 c_{i_0\cdots\widehat{i_j}\cdots i_{q+1}}
 \big|_{U_{i_0\cdots i_{q+1}}}
\]
で定める．
\end{definition}

\begin{proposition}[\v Cech微分の二乗]\label{prop:cech-d2}
$\delta^2=0$ が成立する．
\end{proposition}
\begin{proof}
$\delta(\delta c)$ を展開すると，
添字を二つ取り除く各項が二回現れる．
$i_j$ と $i_k$（$j<k$）を除く二順序は，
一方が符号 $(-1)^{j+k-1}$，
他方が $(-1)^{j+k}$ を持つため打ち消し合う．
制限写像の合成は直接の制限と一致するので，
二項の内容は同じである．
従って全ての項が相殺され，$\delta^2=0$ となる．
\end{proof}

\begin{theorem}[非輪状な被覆による計算]\label{thm:leray}
開被覆 $\mathcal U$ の任意の有限個の元の交わり $V$ について
\[
 H^q(V,\mathcal F)=0\qquad(q\ge1)
\]
が成り立つと仮定する．
このとき
\[
 \check H^q(\mathcal U,\mathcal F)\cong H^q(X,\mathcal F)
\]
が自然に成立する．
\end{theorem}
\begin{proof}
$\mathcal F$ の入射分解
$0\to\mathcal F\to\mathcal I^\bullet$ を取る．
二重複体
\[
 C^{p,q}=
 \prod_{i_0<\cdots<i_p}
 \mathcal I^q(U_{i_0\cdots i_p})
\]
を作る．横微分を\v Cech微分，縦微分を入射分解の
微分とし，全微分を
$d_{\mathrm{tot}}=\delta+(-1)^p d_{\mathrm{vert}}$
と定める．
二つの微分は制限と可換なので，
符号を考慮すれば全微分の二乗は零となる．

まず縦方向のコホモロジーを取る．
各交わり $V$ では入射層の制限が入射層とは限らないが，
入射層はflabbyなので，その制限はflabbyである．
従って制限した入射分解は非輪状分解であり，
その切断複体のコホモロジーは
$H^q(V,\mathcal F)$ を計算する．
仮定によって $q>0$ の項は零で，
$q=0$ の項は $\mathcal F(V)$ である．
よって縦方向からのスペクトル系列は
$q=0$ の行だけに集中し，
全コホモロジーは
$\check H^p(\mathcal U,\mathcal F)$ に等しい．

一方横方向からコホモロジーを取る．
各 $\mathcal I^q$ はflabbyなので，
その任意の局所切断は大域切断へ延長できる．
この性質を用いて，\v Cech複体の正の次数の
コサイクルを順に解消する．
具体的には，被覆に全順序を固定し，
最小の添字を含む局所切断を延長して
零次のcochainを順に調整することで
$\check H^p(\mathcal U,\mathcal I^q)=0$（$p>0$）
を得る．
また $p=0$ のコホモロジーは層の貼り合わせ公理から
$\Gamma(X,\mathcal I^q)$ に等しい．
従って横方向からのスペクトル系列は
$p=0$ 列に集中し，
全コホモロジーは
$H^q(\Gamma(X,\mathcal I^\bullet))
=H^q(X,\mathcal F)$ に等しい．
二つの計算は同じ全複体のコホモロジーを与えるので，
求める自然同型が得られる．
\end{proof}

\section{円周上の局所定数層}

\begin{example}[円周の整数係数コホモロジー]
$S^1$ 上の局所定数整数値関数の層
$\underline{\mathbb Z}$ について
\[
 H^0(S^1,\underline{\mathbb Z})\cong\mathbb Z,\qquad
 H^1(S^1,\underline{\mathbb Z})\cong\mathbb Z,
\]
かつ $q\ge2$ のコホモロジーは零である．
\end{example}
\begin{proof}
円周を三つの開弧 $U_1,U_2,U_3$ で覆い，
各開弧と任意の二つの交わりが連結かつ収縮可能であり，
三つ全ての交わりが空になるように選ぶ．
例えば円周上の中心角が120度ずつ異なり，
長さが120度より大きく240度より小さい開弧を，
三重交わりを生じないように選べる．
収縮可能な開弧上の局所定数層は正の次数の
コホモロジーを持たないので，この被覆は
定理\ref{thm:leray}の仮定を満たす．

\v Cech複体は
\[
 \check C^0=\mathbb Z^3,\qquad
 \check C^1=\mathbb Z^3,\qquad
 \check C^q=0\quad(q\ge2)
\]
となる．
零次微分は
\[
 \delta(a_1,a_2,a_3)
 =(a_2-a_1,a_3-a_1,a_3-a_2)
\]
（適切な辺の順序で記した）である．
その核は三成分が等しい元全体なので
$\ker\delta\cong\mathbb Z$ である．
像は
$(b_{12},b_{13},b_{23})$ について
$b_{13}=b_{12}+b_{23}$ を満たす部分群であり，
その余核は $\mathbb Z$ と同型である．
従って表示したコホモロジー群を得る．
\end{proof}

\begin{remark}
層コホモロジーは，大域切断が右完全でないことを
導来函手として記録する．
大域的な切断の存在障害を測るという見方は，
de Rhamコホモロジーや複素解析における
局所的原始関数の貼り合わせ問題とも関係する．
\end{remark}

\section{到達課題の完全解答}
\noindent\textbf{一点空間上の層コホモロジー．}
\begin{proposition}
一点空間 $X=\{*\}$ 上のAbel群の層 $\mathcal F$
に対して
\[
H^0(X,\mathcal F)=\mathcal F(X),\qquad
H^q(X,\mathcal F)=0\quad(q\ge1)
\]
である．
\end{proposition}
\begin{proof}
一点空間上のAbel群の層の圏は，
全体切断 $\mathcal F\mapsto\mathcal F(X)$ により
Abel群の圏と同値である．
したがって大域切断函手は完全函手である．
層コホモロジーはその右導来函手なので
正次数の導来函手は零である．
零次は定義より大域切断となる．
\end{proof}
\begin{proposition}
一点空間上の短完全列
$0\to\mathcal F'\to\mathcal F\to\mathcal F''\to0$
に大域切断を適用しても短完全である．
\end{proposition}
\begin{proof}
一点空間では全体切断がAbel群の恒等函手と同値なので
完全列を保つ．
\end{proof}
\noindent\textbf{到達確認．}
一点空間上の定数層 $\underline{\mathbb Z}$
の $H^1$ を求めよ．
\par\noindent\textbf{解答．}
上の命題を適用して $H^1(X,\underline{\mathbb Z})=0$．
\end{document}